\documentclass[11pt]{article}
\usepackage[T1]{fontenc}
\usepackage{lmodern}
\usepackage[hmargin=1in,vmargin=0.94in]{geometry}
\usepackage{amsmath,amssymb,amsthm,mathtools,booktabs,microtype}
\usepackage[colorlinks=true,linkcolor=blue,citecolor=blue,urlcolor=blue]{hyperref}
\hypersetup{pdftitle={Uniform intervals for irrational even beta values},pdfauthor={twinheterodox}}
\newtheorem{theorem}{Theorem}[section]
\newtheorem{proposition}[theorem]{Proposition}
\newtheorem{lemma}[theorem]{Lemma}
\newtheorem{corollary}[theorem]{Corollary}
\numberwithin{equation}{section}
\newcommand{\Z}{\mathbb Z}
\newcommand{\Q}{\mathbb Q}
\DeclareMathOperator{\Res}{Res}
\DeclareMathOperator{\lcm}{lcm}
\title{\fontsize{17}{21}\selectfont\mbox{Uniform intervals for irrational even beta values}}
\author{twinheterodox}
\date{18 November 2026}

\begin{document}
\maketitle

\begin{abstract}
For every positive integer \(q\), at least one of the Dirichlet beta values
\(\beta(2q+2),\beta(2q+4),\ldots,\beta(2\lfloor33q/2\rfloor)\) is irrational.
This gives an explicit interval from every even starting index at least four.
For every positive odd \(m\), some even \(s\) with \(m+1\le s\le(299/20)m+C\)
has \(\beta(s)\) irrational, where \(C\) is absolute. This reduces the
previous upper-endpoint multiplier \(18.58\) to \(14.95\) up to a bounded
additive term. For every odd \(m\ge4{,}000{,}000{,}001\), the explicit
endpoint \(s<15m\) suffices.
\end{abstract}

\section{Locating the next irrational value}

Write \(\beta(s)=\sum_{h\ge0}(-1)^h/(2h+1)^s\) for \(s>0\).
Given an even starting index, how far along the list of even beta values
must one go to encounter an irrational number? Knowing that one of
\(\beta(4),\beta(6),\ldots,\beta(32)\) is irrational leaves open the
possibility that \(\beta(4)\) is the only irrational member of that
list. To answer the question starting at \(\beta(10)\), or at a
starting index chosen arbitrarily far out, one needs control of both
ends of the interval.

\begin{theorem}\label{th:all}
For every integer \(q\ge1\), some even integer \(s\) satisfies
\[
             2q+2\le s\le 2\lfloor 33q/2\rfloor,
             \qquad \beta(s)\ \text{is irrational}.
\]
\end{theorem}

The first three upper endpoints are \(32,66,98\), at starting indices
\(4,6,8\), respectively. The first is known from \cite{fifteen}, and
the next two from \cite{two}. Starting at \(10\) gives the list
\(\beta(10),\beta(12),\ldots,\beta(132)\). The theorem gives an endpoint
for each subsequent starting index as well.

\begin{theorem}\label{th:slope}
There is an absolute constant \(C\) such that for every positive odd
integer \(m\) there is an even \(s\) with
\[
         m+1\le s\le \frac{299}{20}m+C,\qquad
         \beta(s)\ \text{irrational}.
\]
For every odd \(m\ge4{,}000{,}000{,}001\), one can take \(m+1\le s<15m\).
\end{theorem}

The earlier uniform result gives the explicit upper endpoint
\(m-1+2\lceil(879m-50)/100\rceil<18.58m\) for every positive odd \(m\)
\cite[Proposition B.1]{shifted}. Theorem~\ref{th:all} gives an
asymptotic multiplier \(16.5\) on writing \(m=2q+1\).
The first assertion of Theorem~\ref{th:slope} improves this to \(14.95\),
and the second gives a specified range for a multiplier of \(15\).

Ball and Rivoal used reflection in a rational function to isolate odd
zeta values \cite[Section~1]{ballrivoal}. For even beta values,
Rivoal and Zudilin used half-integer poles and an alternating sum
\cite[Lemma~2]{rzbeta}. Rivoal used a second derivative to remove
\(\zeta(3)\) from odd-zeta forms \cite[Section~2, Lemma~1]{rivoal}.
Krisman adapted that step to remove \(\beta(2)\), giving the integer
zeros multiplicity three so that they survive differentiation
\cite[Section~1.2]{krisman}. We use the order-\(31\) denominator
refinement of \cite{fifteen}. Taking a power before differentiating
makes the starting index adjustable \cite{powered,two}.
Section~\ref{sec:fixed} treats two choices of derivative order. In each
case irrationality in the interval follows from the nonvanishing of
a particular rational number, which we call its \emph{seed}.

Two observations let finitely many computations control infinitely many
intervals. First, a successful power can serve several starting points.
For example, the ninth power with twenty derivatives gives the interval
from \(22\) to \(298\). It fits inside the interval from \(22\) to
\(330\) requested at \(q=10\). Its multiples have lower and upper
indices \(20j+2\) and \(299j\), which explains the ratio \(299/20\).
Second, the normalized seeds are coefficients of an algebraic function.
Its behavior at infinity forces bounded gaps between nonzero
coefficients. Reduction modulo primes lets a single rational seed
certify almost every prime exponent. We use bounded gaps for the
sharpest slope and primes for the explicit ranges.

\section{From a rational seed to an irrational value}\label{sec:fixed}

The denominators consist of blocks of consecutive poles of different
lengths. Zudilin used such blocks for even beta values
\cite[Section~3]{zudbeta}. Lai and Zhou changed their lengths and
factorial normalization to reach the interval from \(\beta(2)\) to
\(\beta(10)\) \cite[Theorem~6.1]{laizhou}. Here we take the \(31\)
blocks specified in \cite[Section~2]{fifteen}.

Let \(96\mid n\), \(n>0\). For \(0\le h\le30\) put
\(I_h=[hn/96,(1-h/96)n]\cap\Z\), and let \(B_a(n)\) count the
sets \(I_h\) containing \(a\). With
\(\Lambda_n=[-7n/4-1,3n/4]\cap\Z\), the kernel is
\[
 R_n(t)=\frac{(2t+n+1)\prod_{v\in\Lambda_n}(t-v)^3}
                {\prod_{a=0}^{n}(t+a+\tfrac12)^{B_a(n)}} .
\]
Its symmetry is \(R_n(-n-1-t)=R_n(t)\). For fixed \(A\ge1\) and even
\(d\ge2\) we use the form
\[
       L_n(A,d)=\sum_{t\ge0}(-1)^t (R_n^A)^{(d)}(t).
\]
The two pairs throughout the proof are
\begin{equation}\label{eq:pairs}
                  (A,d)=(j,2j)\quad\text{or}\quad(9j,20j),
                  \qquad j\in\Z_{\ge1}.
\end{equation}
At a pole of order \(i\), differentiation raises the order to \(i+d\).
Reflection cancels the odd \(i\) in the beta coefficients.
Consequently the even beta indices range from \(d+2\) to at most
\(31A+d\). In both pairs \(d<3A\), so the derivative still vanishes
at each integer zero. This inequality will also be needed to clear
denominators.

The seed is computed at just 97 pole addresses. Set
\[
 u_a=96-2a,\quad B_a=B_a(96)=1+\min(a,96-a,30)\quad(0\le a\le96),
 \qquad \sigma_a(s)=\sum_{h=0}^{a-1}\frac{(-1)^h}{(2h+1)^s}.
\]
Define the even rational function \(f\) and its polar coefficients by
\begin{equation}\label{eq:f}
 f(z)=\frac{z\prod_{\substack{|v|\le239\\v\ {\rm odd}}}(z-v)^3}
                     {\prod_{a=0}^{96}(z-u_a)^{B_a}},
 \qquad f(z)^A=\sum_{a=0}^{96}\sum_{i=1}^{A B_a}\frac{c_{A,a,i}}{(z-u_a)^i}.
\end{equation}
The cutoff \(239\) comes from the zeros divisible by the dilation prime,
as in \cite[Section 5]{fifteen}. A nominal coefficient can vanish at the
central pole because the numerator contains \(z\). The normalized seed is
\begin{equation}\label{eq:seed}
 b(A,d)=-\sum_{a=0}^{96}(-1)^a\sum_{i=1}^{A B_a}
       \binom{i+d-1}{d}\,c_{A,a,i}\,\sigma_a(i+d).
\end{equation}
Multiplying by \(d!\) restores the rising factorial
\((i)_d=i(i+1)\cdots(i+d-1)\) from differentiation.

\begin{proposition}\label{prop:fixed}
For either pair in \eqref{eq:pairs}, if \(b(A,d)\ne0\), some even
\(s\in[d+2,31A+d]\) has \(\beta(s)\) irrational.
\end{proposition}

Zudilin bounded the polar coefficients by differentiating the
individual factors \cite[proof of Lemma~4]{zudbeta}. Krisman's
valuation bound for the series at each pole \cite[Lemma~4.3]{krisman}
is the starting point for the isolated-pole estimate for these \(31\)
blocks. We use that estimate and the real phase bounds from
\cite[Sections~3--5 and Appendix~A]{fifteen}. The case \(d=2A\)
after powering is proved in \cite[Sections~3--5]{two}.

\subsection{The cost of the extra derivatives}

Write \(D_n=\prod_{p\le n}p^{e_p}\) for the integer of
\cite[Section 3 and Appendix A]{fifteen}, and put
\[
 \lambda_p=\lfloor\log_p(7n)\rfloor\quad(p\text{ odd}),\qquad
 H_n=\prod_{\substack{p\le n\\p\ {\rm odd\ prime}}}p^{\lambda_p}.
\]
We claim that the rational coefficients of \(L_n(A,d)\), including
its constant term, are integers after multiplication by
\begin{equation}\label{eq:clear}
                          D_n^A H_n^{\,d-2A}.
\end{equation}
To see where the extra cost enters, write \(a_{n,A;a,i}\) for
the coefficient of \((t+a+\tfrac12)^{-i}\) in \(R_n(t)^A\).
Partial fractions express the form as
\[
 L_n(A,d)=\sum_{a,i}(-1)^a 2^{i+d}(i)_d\,a_{n,A;a,i}
                       \{\beta(i+d)-\sigma_a(i+d)\}.
\]
The coefficients at \(a\) and \(n-a\) differ by the factor \((-1)^i\).
This proves the parity and endpoint assertions above. All the
shifts and differentiations involve absolutely convergent series.

At an odd \(p\) with \(p^2\le7n\), write
\[
 \nu_p(a)=3\sum_{v\in\Lambda_n}v_p(2a+1+2v)
             -\sum_{b\ne a} B_b(n)v_p(b-a).
\]
By \cite[Appendix A]{fifteen}, the degree-\(h\) coefficient of the
series obtained by removing the nominal pole of \(R_n\)
has valuation at least \(\nu_p(a)-h\lambda_p\), and
\(e_p\ge-\nu_p(a)+33\lambda_p\). Multiplication of series gives
\[
  v_p(a_{n,A;a,i})\ge A\nu_p(a)-\lambda_p(AB_a(n)-i).
\]
A harmonic denominator costs at most \((i+d)\lambda_p\). Since
\(B_a(n)\le31\), the combined loss is at most \(Ae_p+(d-2A)\lambda_p\).
Such primes are at most \(n\). At \(2\) all harmonic denominators are
odd, so the inherited local bound and \(D_n^A\) already suffice.

For odd \(p\le n\) with \(p^2>7n\), addresses are grouped by their
residue modulo \(p\). In a repeated class the local estimate is
\(v_p(a_{n,A;a,i})\ge i-AF_p(a)\), where the local pole excess is
\[
 F_p(a)=\sum_{\substack{0\le b\le n\\b\equiv a\bmod p}} B_b(n)
        -3\#\{v\in\Lambda_n:p\mid 2a+1+2v\}-\mathbf 1_{p\mid n-2a}.
\]
If \(M_p\) is its maximum over the repeated classes, then
\(e_p=\max(0,M_p+2)\) \cite[Section 3]{fifteen}.
An evaluation denominator divisible by \(p\) costs \(i+d\) powers,
so the total valuation is at least \(-AM_p-d\). This is covered by
\(Ae_p+d-2A\), and here \(\lambda_p=1\). In a class with only one
address the local coefficients are \(p\)-integral.

For the last assertion to help with the harmonic denominators, group
the constant term by its evaluation point \(-\ell\), \(1\le\ell\le n\).
The term at this point uses addresses \(a\ge\ell\) with denominators
\((a-\ell+\tfrac12)^{i+d}\). Summing over \emph{all} addresses gives
\((R_n^A)^{(d)}(-\ell)=0\), since \(d<3A\). Thus the sum over
\(a\ge\ell\) equals the negative of the sum over \(a<\ell\).
Zudilin used zeros in the pole interval to shift the summation limit
in the rational term \cite[proofs of Lemmas~1 and~5]{zudbeta}.
Krisman obtained the two complementary sums after two derivatives
\cite[Lemma~4.6]{krisman}. Here \(d<3A\) preserves the same choice
for the powered kernel.
Only the class \(2(a-\ell)+1\equiv0\pmod p\) can have a bad
denominator. If it has just one address, choose the side which omits
that address. All other terms then have unit denominators and the
stated local coefficient bounds apply. For \(p>n\) every class has at
most one address. This finishes the proof of \eqref{eq:clear}.

First powers contribute \(\sum_{2<p\le n}\log p\) to \(\log H_n\),
and higher powers contribute \(O(\sqrt n\log n)\). Thus
\(\log H_n=(1+o(1))n\) by the prime number theorem. We compare this cost with the decay of the
same form. The first nonzero summand can occur only at \(t>3n/4\).
On these integers logarithmic differentiation gives
\(\lvert(R_n^A)^{(d)}(t)\rvert\le C_{A,d}n^d R_n(t)^A\), since all the
linear factors have distance at least one from their zeros or poles
and there are \(O(n)\) factors counted with multiplicity.
The Stirling and prime-density estimates of \cite[Section 4]{fifteen}
give the following upper bound with \(\delta=41/5+15\log2-18.83<-0.23\):
\begin{equation}\label{eq:decay}
 \limsup_{\substack{n\to\infty\\96\mid n}}\frac1n
 \log\bigl|D_n^A H_n^{d-2A}L_n(A,d)\bigr|
                 \ \le\ A\delta+(d-2A)<0 .
\end{equation}
Indeed the \(G\log(n!)\) term in \(\log D_n\), \(G=221/16\),
cancels the corresponding term from the kernel. The remaining
denominator cost is at most \(41/5+15\log2+o(1)\) per \(n\), and the
real phase is at most \(-18.83\). Beyond \(t=3n\) its additional
bound \((3n/t)^{3An}\) absorbs the polynomial factors from Stirling's
formula and differentiation. All these estimates take \(A,d\) fixed.
For the ninth-power family the two extra derivatives per \(j\) cost
\(2j\), leaving \(9j\delta+2j<-0.07j\).

\subsection{How the seed prevents vanishing}

Lai and Sprang's real auxiliary-prime criterion singles out the
constant coefficient by its strictly smaller valuation at primes
tending to infinity \cite[Lemma~2.3]{laisprang}. For this kernel,
the dilation developed in \cite[Section~4]{powered} and used in
\cite[Section~5]{fifteen,two} produces that separation.
Set \(n=96P\) with \(P>337\) prime and
use the centered variable \(x=2t+n+1\). In this variable the kernel,
up to the factor \(2^{Gn+25}\), is
\[
 F_n(x)=\frac{x\prod_{v\ {\rm odd},\,|v|\le5n/2+1}(x-v)^3}
                    {\prod_{a=0}^n(x-n+2a)^{B_a(n)}} .
\]
Write \(\gamma_{n,A;a,i}\) for the polar coefficient of \(F_n^A\)
at address \(a\). After removal of the common nonzero factor
\(2^{A(Gn+25)+d}\), the form \(L_n(A,d)\) is
\[
 \sum_{a,i}(-1)^a(i)_d\,\gamma_{n,A;a,i}
                          \{\beta(i+d)-\sigma_a(i+d)\}.
\]
The dilation lemma of \cite[Section 5]{two} applies directly to the
undifferentiated power. It gives
\[
 P^{1356A-i}\gamma_{96P,A;Pa,i}\equiv(-1)^A c_{A,a,i}\pmod P,\qquad
 v_P(\gamma_{96P,A;b,i})\ge i-1326A\quad(P\nmid b).
\]
The first congruence includes integrality at \(P\).
Multiply the centered form by \(P^{1356A+d}\). At a selected address
\(Pa\), the order-\(i\) term has \(i+d\) powers of \(P\) available.
Its beta coefficient vanishes modulo \(P\). In its harmonic correction
only the terms \(2h+1=P(2v+1)\) survive, giving
\[
        P^{i+d}\sigma_{Pa}(i+d)\equiv
        (-1)^{(P-1)/2}\sigma_a(i+d)\pmod P .
\]
For \(b\) outside the selected addresses, each harmonic denominator
contains at most one factor \(P\), since \(2n<P^2\).
Even after this loss at least \(30A\) further powers remain. Thus
the scaled rational constant reduces to
\((-1)^{A+(P-1)/2}d!\,b(A,d)\), and every scaled beta coefficient
is divisible by \(P\).

\begin{proof}[Conclusion of Proposition~\ref{prop:fixed}]
Suppose that the beta values in the indicated interval are rational
with common denominator \(T\). If \(b(A,d)\ne0\), every sufficiently
large prime \(P\) avoiding \(T\) and the numerator and denominator of
\(d!\,b(A,d)\) gives a nonzero form at \(n=96P\). Indeed its scaled
constant is a unit at \(P\) while every other scaled term vanishes
modulo \(P\). By \eqref{eq:clear} and \eqref{eq:decay}, the numbers
\(T D_n^A H_n^{d-2A}L_n(A,d)\) along that subsequence are nonzero
integers tending to zero. This contradiction proves the proposition.
\end{proof}

\section{Why the successful powers have bounded gaps}\label{sec:gaps}

We use the same argument for both derivative orders. Fix
\((k,r)=(1,2)\) or \((9,20)\), and set \(b_j=b(kj,rj)\) for \(j\ge1\)
and \(b_0=0\). The first term is nonzero in each case. For the first
sequence \(2!\,b(1,2)\equiv3738\pmod{10007}\)
\cite[Section 5]{fifteen}. For the second the following finite-field
calculations use the normalization in \eqref{eq:seed}.
\begin{center}
\begin{tabular}{@{}rrr@{}}\toprule
Prime & \(b(9,20)\) & \(20!\,b(9,20)\)\\\midrule
\(10007\)&\(1491\)&\(3058\)\\
\(10009\)&\(7947\)&\(4140\)\\\bottomrule
\end{tabular}
\end{center}
At each prime we compute the \(9\sum_a B_a=18{,}693\) nominal polar
coefficients both by multiplying truncated local factors and by the
logarithmic derivative recurrence of \cite[Section 6]{two}. We then
sum \eqref{eq:seed} in its original order and in the order used
in \eqref{eq:residues} below. All inversions occur below the
characteristic or at unit pole differences. The program
\texttt{checks/derivative\_seed.py} in the accompanying
\texttt{uniform-even-beta-intervals} source directory reproduces
the table with the option \texttt{--cross-check}. The same directory
convention is used for the other programs below. The original kernel
certificates are in the project repository \cite{repo}.

\begin{theorem}\label{th:seeds}
For each of the two sequences there is an integer \(L\ge1\) such that
every \(L\) consecutive positive indices contain a \(j\) with \(b_j\ne0\).
Moreover, for every prime \(p>192\) and every \(j\ge1\),
\[
                         b_{pj}\equiv b_j\pmod p
                         \qquad\text{in }\Z_{(p)}.
\]
In particular, if \(b_j\ne0\), then for all but finitely many primes
\(p\) all coefficients \(b_{jp^e}\), \(e\ge0\), are nonzero.
\end{theorem}

\subsection{A residue representation with a fixed rational function}

The difficulty with \(b_j\) in \eqref{eq:seed} is that both the order of
the pole and the order of differentiation grow with \(j\).
The normalization by \((rj)!\) allows us to put both changes in
one ordinary power. For \(1\le\ell\le96\) define
\[
                    y_\ell=97-2\ell,\qquad
                    g_\ell(z)=\frac{f(z)^k}{(y_\ell-z)^r}.
\]
In \eqref{eq:seed}, change the harmonic summation index \(h\) to
\(\ell=a-h\). Then \(y_\ell-u_a=2h+1\) and
\((-1)^{a+h}=(-1)^\ell\). Expanding \((y_\ell-z)^{-rj-1}\) around \(u_a\)
gives, as its coefficient of \((z-u_a)^{i-1}\),
\(\binom{i+rj-1}{rj}(y_\ell-u_a)^{-i-rj}\). Thus
\begin{equation}\label{eq:residues}
 b_j=-\sum_{\ell=1}^{96}(-1)^\ell
                 \sum_{a=\ell}^{96}
                 \Res_{z=u_a}\frac{g_\ell(z)^j}{y_\ell-z}\,dz .
\end{equation}
This holds also at \(j=0\), when every residue is zero. Each \(y_\ell\)
is a triple zero of \(f\). Thus \(g_\ell\) has a zero there of order
\(3k-r>0\). Its other finite poles are among the \(u_a\). At infinity
\(f\) has a zero of order \(1356\), by counting the factors of
\eqref{eq:f}. Consequently \(g_\ell\) has a zero at infinity of
order \(1356k+r\).

\subsection{Algebraicity and propagation of zeros}

Consider \(B(T)=\sum_{j\ge0}b_jT^j\). For each \(\ell\), take disjoint
small circles around the \(u_a\) occurring in \eqref{eq:residues},
with \(y_\ell\) outside. For \(|T|\) sufficiently small the geometric
series converges uniformly on all the circles. Summing under their
integrals replaces the differential in \eqref{eq:residues} by
\[
                    \frac{dz}{(y_\ell-z)(1-Tg_\ell(z))}.
\]
For a generic nonzero \(T\) near zero the original poles \(u_a\)
are regular points of this differential. The poles now enclosed
are the simple roots \(\xi\) of \(g_\ell(\xi)=1/T\) tending to the
chosen centers. Indeed, writing \(g_\ell=P/Q\) in lowest terms,
Rouch\'e's theorem on each circle counts the roots of \(Q-TP\)
there with their multiplicities. The critical parameter values form
a finite set. The residue at such a simple root is
\begin{equation}\label{eq:rootres}
 h_\ell(\xi)=\frac{g_\ell(\xi)}
                      {(\xi-y_\ell)g_\ell'(\xi)}.
\end{equation}
Thus \(B(T)\) is a finite signed sum of algebraic functions over
\(\Q(T)\), and is itself algebraic.

To show that the series has infinitely many nonzero coefficients,
follow the roots to large \(|T|\).
They approach zeros of \(g_\ell\) on the Riemann sphere. At a finite
zero \(v\) of order \(e\), we have \(g_\ell/g_\ell'=(\xi-v)/e+
O((\xi-v)^2)\). The residue \eqref{eq:rootres} tends to \(1/e\) if
\(v=y_\ell\), and to zero for every other finite zero. At infinity
it tends to \(-1/(1356k+r)\). Since every root lies near one of
these zeros for sufficiently large \(|T|\), all branches of the sum
are bounded there. A polynomial \(B(T)\) would therefore be constant.
Its linear coefficient \(b_1\) is nonzero, so its coefficients cannot
eventually vanish.

The coefficients of an algebraic power series satisfy a recurrence with
polynomial coefficients \cite[Theorems~1.5 and~2.1]{stanley}. We recall
the argument to see why a run of zero coefficients propagates.
The finite extension \(\Q(T,B)\) is closed under differentiation.
For a minimal polynomial \(A(T,X)\), implicit differentiation gives
\(B'=-A_T(T,B)/A_X(T,B)\), since characteristic zero gives separability.
The derivatives of \(B\) are therefore linearly dependent over
\(\Q(T)\). Clearing denominators and comparing coefficients gives a
recurrence, after translating the index by a fixed amount,
\begin{equation}\label{eq:rec}
       \sum_{i=0}^{L}P_i(j)b_{j+i}=0\quad(j\text{ sufficiently large}),
       \qquad P_i\in\Q[j],\quad P_L\ne0.
\end{equation}
Indeed, the terms \(T^v B^{(u)}\) of the differential equation
contribute shift \(u-v\). Among terms of any fixed shift, the term
of highest derivative order has strictly highest polynomial degree
in \(j\). Thus some coefficient polynomial survives. The case \(L=0\)
would force eventual vanishing. Beyond the last integer root of
\(P_L\), a run of \(L\) zeros propagates forwards by \eqref{eq:rec}
and would also force eventual vanishing. Enlarging \(L\) covers the
finite initial segment. This proves the bounded-gap assertion.

\subsection{Prime exponents}

Fix \(p>192\). Distinct pole positions differ by at most \(192\), and
\(y_\ell-u_a\) is a nonzero odd integer of absolute value at most
\(191\). Consequently the Laurent expansion of \(g_\ell^j\) at
\(w=z-u_a\) has \(p\)-integral coefficients. Write its negative
part as \(\sum_{i\ge1}a_{-i}w^{-i}\), a finite sum. In characteristic
\(p\), taking the \(p\)th power sends \(w\) to \(w^p\) and fixes the
reduced coefficients. With \(v=y_\ell-u_a\), the two residues after
division by \(y_\ell-z=v-w\) reduce respectively to
\[
                  \sum_{i\ge1}a_{-i}v^{-i},
                  \qquad \sum_{i\ge1}a_{-i}v^{-pi}.
\]
They agree by Fermat's theorem. Summing \eqref{eq:residues} proves the
congruence and hence Theorem~\ref{th:seeds}. For any fixed nonzero
rational \(b_j\), every sufficiently large prime is a unit prime for
\(b_j\), so the congruence can be iterated at that same prime for all
\(e\). The corresponding unnormalized seed \((rjp^e)!\,b_{jp^e}\)
is nonzero over \(\Q\).

\begin{proof}[Proof of the first assertion of Theorem~\ref{th:slope}]
Take \(L\) for the family \((9,20)\). Given a positive odd \(m\), set
\(J=\max(1,\lceil(m-1)/20\rceil)\). There is a successful
\(J\le j\le J+L-1\). Proposition~\ref{prop:fixed} gives an irrational
beta value at an even index between \(20j+2\) and \(299j\).
The first endpoint is at least \(m+1\). We can take \(C=299L\), since
\[
            299j\le \frac{299}{20}(m-1)+299L .\qedhere
\]
\end{proof}

For the original family \(k_j=(2j)!\,b(j,2j)\), Theorem~\ref{th:seeds}
also supplies bounded gaps and all powers of every sufficiently large
prime. Whether \(k_j\ne0\) at every sufficiently large \(j\) remains
open. Covering a prescribed interval leaves us free to use any
successful pair whose two endpoints fit inside it.

\section{An effective supply of prime exponents}\label{sec:effective}

The bounded-gap argument gives no numerical value for \(L\). We can
obtain explicit endpoints by choosing the exponent \(j\) to be a
prime. Write \(B_*=b(9,20)\). For every prime \(p>192\) at which
\(B_*\) is a unit, Theorem~\ref{th:seeds} supplies the interval
\([20p+2,299p]\). To find such a prime in a prescribed range, we
compare the logarithmic weight of \emph{all} possible exceptional
primes with the weight of the primes in the range.

\begin{lemma}\label{lem:height}
There is a nonzero integer \(M\) with \(\log|M|<180{,}000\) divisible
by every prime \(p>192\) at which \(B_*\) is not a unit.
\end{lemma}
\begin{proof}
Put \(D=\lcm(1,\dots,192)\) and
\[
 D_a=\prod_{b\ne a}|u_a-u_b|^{9B_b},\qquad
 Q=D^{299}\lcm_{0\le a\le96}D_a .
\]
Then \(M=QB_*\) is a nonzero integer. Indeed, expand
\(w^{9B_a}f(u_a+w)^9\). In degree \(h\), the other pole factors
introduce denominators dividing \(D_aD^h\), by the integral binomial
expansion of each inverse linear factor. For the pole coefficient of
order \(i\) we take \(h=9B_a-i\). Each accompanying harmonic
denominator divides \(D^{i+20}\), and \(h+i+20\le299\). This proves
integrality. All prime divisors of \(Q\) are at most \(192\).

On \(|w|=1\), every other pole factor has modulus at least one.
The bounds \(|u_a+w|\le97\) and \(|u_a+w-v|\le336\) show by Cauchy's
coefficient estimate that \(|c_{9,a,i}|\le97^9 336^{6480}\).
Each of the \(18{,}693\) terms has harmonic sum of magnitude at most
\(96\) and binomial factor at most \(2^{299}\). Hence
\[
      1\le |M|\le Q\cdot18{,}693\cdot96\cdot2^{299}97^9 336^{6480}.
\]
This upper bound is an explicitly specified integer of bit length
\(256{,}201\), as reproduced by \texttt{checks/explicit\_threshold.py}.
Since \(\log2<7/10\), its logarithm is less than \(180{,}000\).
\end{proof}

\begin{proposition}\label{prop:bands}
For \(x\ge4{,}000{,}000\), the interval \((x,11x/10]\) contains a
prime \(p\) with \(b(9p,20p)\ne0\). For \(x\ge200{,}000{,}000\)
the same holds in \((x,1003x/1000]\).
\end{proposition}
\begin{proof}
Let \(\vartheta(x)=\sum_{p\le x}\log p\). Dusart's estimate
\cite[Theorem 5.2, table on p.~4]{dusart} gives
\[
       |\vartheta(x)-x|\le\frac{x}{5\log^2 x}
                  \qquad(x>3{,}594{,}641).
\]
If \(x\ge x_0\) and \(\log x_0>t\), subtracting the two estimates
at \(x\) and \((1+u)x\) gives a strict lower bound
\(x[u-(2+u)/(5t^2)]\) for the logarithmic weight in the interval.
In the two cases its value at \(x_0\) is as follows.
\begin{center}
\begin{tabular}{@{}rrrr@{}}\toprule
\(x_0\) & \(t\) & \(u\) & \(x_0[u-(2+u)/(5t^2)]\)\\\midrule
\(4{,}000{,}000\)&14&\(1/10\)&\(2{,}740{,}000/7\)\\
\(200{,}000{,}000\)&18&\(3/1000\)&\(28{,}570{,}000/81\)\\\bottomrule
\end{tabular}
\end{center}
The comparisons \((11/4)^t<x_0\) and \(e<11/4\) verify the log
hypotheses. Both lower bounds exceed \(180{,}000\). By
Lemma~\ref{lem:height}, some prime \(p\) in each such interval does
not divide \(M\). At that prime \(b(9p,20p)\equiv B_*\ne0\pmod p\).
\end{proof}

\begin{proof}[Explicit endpoints for large starting indices]
For \(q\ge40{,}000{,}000\) take \(x=q/10\) in the first prime interval.
Its successful exponent \(p\) satisfies
\[
       20p+2>2q+2,\qquad 299p\le(299\cdot11/100)q<33q .
\]
The even index supplied by Proposition~\ref{prop:fixed} is at most
\(2\lfloor33q/2\rfloor\). For odd \(m\ge4{,}000{,}000{,}001\), take
\(x=(m-1)/20\) in the second interval. Now \(20p+2>m+1\) and
\(299p\le(299\cdot1003/20000)(m-1)<15m\). This proves the second
assertion of Theorem~\ref{th:slope}. In both applications the exponent
prime \(p\) has been fixed before the dilation prime \(P\) in
Proposition~\ref{prop:fixed} tends to infinity.
\end{proof}

\section{Covering the finite initial range}\label{sec:finite}

It remains to prove Theorem~\ref{th:all} below \(q=40{,}000{,}000\).
One successful pair \((9j,20j)\) works for every \(q\) satisfying
\begin{equation}\label{eq:cover}
                 \left\lceil\frac{299j}{33}\right\rceil
                          \le q\le 10j .
\end{equation}
Indeed, these inequalities place \([20j+2,299j]\) inside
\([2q+2,33q]\), and rounding both upper endpoints down to even integers
preserves containment. This is why forty million starting points can be
covered by a short list of certificates. We describe the selection
rules and the exact arithmetic verifying them. All step-by-step
residues and endpoints accompany the programs in \texttt{checks/}.

\paragraph{The middle range.}
Start with the prime \(193\). Given the last prime \(p\), append the
largest prime at most \(\lfloor27p/25\rfloor\), stopping when
\(10p\ge40{,}000{,}000\). For each entry test the same rational
number \(B_*=b(9,20)\) modulo that prime. The computation
\texttt{prime\_bridge.py} verifies primality by trial division and
strict increase. It evaluates \(B_*\) modulo the product of the
distinct primes and the two reference primes \(10007,10009\).
The denominator factors inverted in the truncated products are units
in every component, since all primes exceed \(192\). The binomial
coefficients in \eqref{eq:seed} are formed as integers before
reduction. Thus the product calculation gives the individual residues
exactly, even for primes smaller than \(299\).

\begin{center}
\begin{tabular}{@{}lrrr@{}}\toprule
& Exponent \(p\)& \(B_*\bmod p\)& Range in \(q\) from \eqref{eq:cover}\\\midrule
First & \(193\)&\(27\)& \(1749\) to \(1930\)\\
Last & \(4{,}047{,}157\)&\(660{,}746\)&
 \(36{,}669{,}696\) to \(40{,}471{,}570\)\\\bottomrule
\end{tabular}
\end{center}
All \(135\) residues are nonzero. The recorded neighboring intervals
satisfy \(\mathrm{left}_{i+1}\le\mathrm{right}_i+1\).
Theorem~\ref{th:seeds} therefore gives complete coverage from \(1749\)
through \(40{,}471{,}570\).

\paragraph{The first starting points.}
Begin with \(q=1\), and always take the first uncovered \(q\).
Set \(j=\lfloor33q/299\rfloor\). If \(j\ge1\) and \(q\le10j\),
test \(b(9j,20j)\) and cover through \(10j\). Otherwise test
\(b(q,2q)\) and cover \(q\) itself. Stop once the next uncovered
\(q\) is at least \(1749\). The resulting list consists of \(89\)
seeds, \(45\) of the original type and \(44\) of the ninth-power
type. Every residue is nonzero modulo \(65537\), and the intervals
cover \(1\) through \(1890\). The largest \(31A+d\) used is \(56{,}511\).

To certify these Taylor coefficients, first extract the central zero.
At \(u_a\) the remaining products \(P,Q\) of linear factors from
\eqref{eq:f} have unit constant terms modulo \(65537\). Thus
\[
 w^{B_a}f(u_a+w)=w^\epsilon P(w)/Q(w),
 \qquad \epsilon=\mathbf 1_{u_a=0}.
\]
For \(A\) in the list put \(l=A(B_a-\epsilon)\). The
coefficients of \(Y=(P/Q)^A\) through degree \(l-1\) are uniquely
determined in the field by
\[
  Y(0)=(P(0)/Q(0))^A,\qquad
  P Q\,Y'=A(P'Q-PQ')Y \pmod{w^{\,l-1}} .
\]
In degree \(h-1\) the new coefficient of \(Y\) is multiplied by
\(hP(0)Q(0)\), a unit for \(1\le h<l<65537\).
The polar coefficient of order \(i\) is \([w^{\,l-i}]Y\) for
\(1\le i\le l\), and is zero for \(l<i\le AB_a\).
The program \texttt{initial\_intervals.py} verifies these identities
on \(21{,}187{,}413\) coefficients before using the symmetry
\(c_{A,96-a,i}=(-1)^i c_{A,a,i}\) for the remaining addresses.
It evaluates \eqref{eq:seed} with the finite harmonic sums.

As a check on the computation, \texttt{interval\_replay.py} computes
all \(89\) seeds modulo a second prime, \(100003\), expanding every
contributing pole directly and summing the harmonic correction in
reverse order. All \(89\) are again nonzero. It also checks the \(135\)
prime residues one field at a time. The accompanying tables can be
recomputed together with \texttt{python3 -B checks/verify.py} from the
\texttt{uniform-even-beta-intervals} directory, using Python and
\texttt{python-flint}.

\begin{proof}[Conclusion of Theorem~\ref{th:all}]
The two finite ranges \([1,1890]\) and \([1749,40{,}471{,}570]\)
overlap. The latter overlaps the range \(q\ge40{,}000{,}000\)
proved in Section~\ref{sec:effective}. Hence every positive integer
\(q\) has a successful pair in \eqref{eq:pairs} whose interval fits
inside the prescribed one. Proposition~\ref{prop:fixed} supplies an
irrational beta value there.
\end{proof}

\paragraph{Verification.}
Independent AI reviewers checked the mathematical arguments in the
17 November 2026 proof records and reproduced the seed, coverage and
threshold checks. The effective-threshold argument also received a
separate review of its exposition.

\begin{thebibliography}{99}
\small
\bibitem{shifted}
twinheterodox, \emph{Irrational even beta values in shifted intervals},
2026, \href{https://hexagonmath.org/2611.00160v2}{Hexagon 2611.00160v2}.
\bibitem{fifteen}
twinheterodox, \emph{At least one of \(\beta(4),\beta(6),\ldots,\beta(32)\)
is irrational}, 2026,
\href{https://hexagonmath.org/2611.00346v1}{Hexagon 2611.00346v1}.
\bibitem{two}
twinheterodox, \emph{Two shorter intervals of even Dirichlet beta values},
2026, \href{https://hexagonmath.org/2611.00601v1}{Hexagon 2611.00601v1}
(\href{https://github.com/twinheterodox/beta4/tree/main/beta-intervals}{manuscript and certificates}).
\bibitem{ballrivoal}
K.~Ball and T.~Rivoal, \emph{Irrationalit\'e d'une infinit\'e de valeurs de la
fonction z\^eta aux entiers impairs}, Invent. Math. \textbf{146} (2001),
193--207, \href{https://doi.org/10.1007/s002220100168}{doi:10.1007/s002220100168}.
\bibitem{rzbeta}
T.~Rivoal and W.~Zudilin, \emph{Diophantine properties of numbers related to
Catalan's constant}, Math. Ann. \textbf{326} (2003), 705--721,
\href{https://doi.org/10.1007/s00208-003-0420-2}{doi:10.1007/s00208-003-0420-2}.
\bibitem{krisman}
S.~Krisman, \emph{One of \(\beta(4),\beta(6),\ldots,\beta(36)\) is irrational},
2026, \href{https://hexagonmath.org/2610.00285v2}{Hexagon 2610.00285v2}.
\bibitem{rivoal}
T.~Rivoal, \emph{Irrationalit\'e d'au moins un des neuf nombres
\(\zeta(5),\zeta(7),\ldots,\zeta(21)\)}, Acta Arith. \textbf{103}
(2002), 157--167, \href{https://doi.org/10.4064/aa103-2-5}{doi:10.4064/aa103-2-5}.
\bibitem{powered}
twinheterodox, \emph{Irrational even beta values from powered hypergeometric
kernels}, 2026, \href{https://hexagonmath.org/2611.00336v1}{Hexagon 2611.00336v1}.
\bibitem{zudbeta}
W.~Zudilin, \emph{Arithmetic of Catalan's constant and its relatives},
Abh. Math. Semin. Univ. Hambg. \textbf{89} (2019), 45--53,
\href{https://doi.org/10.1007/s12188-019-00203-w}{doi:10.1007/s12188-019-00203-w}.
\bibitem{laizhou}
L.~Lai and L.~Zhou, \emph{At least two of \(\zeta(5),\zeta(7),\ldots,\zeta(35)\)
are irrational}, Publ. Math. Debrecen \textbf{101} (2022), 353--372,
\href{https://doi.org/10.5486/PMD.2022.9252}{doi:10.5486/PMD.2022.9252}.
\bibitem{laisprang}
L.~Lai and J.~Sprang, \emph{Many \(p\)-adic odd zeta values are irrational},
2025, \href{https://arxiv.org/abs/2306.10393v2}{arXiv:2306.10393v2}.
\bibitem{stanley}
R.~P.~Stanley, \emph{Differentiably finite power series}, European J. Combin.
\textbf{1} (1980), 175--188,
\href{https://doi.org/10.1016/S0195-6698(80)80051-5}{doi:10.1016/S0195-6698(80)80051-5}.
\bibitem{dusart}
P.~Dusart, \emph{Estimates of some functions over primes without R.H.},
2010, \href{https://arxiv.org/abs/1002.0442v1}{arXiv:1002.0442v1}.
\bibitem{repo}
twinheterodox, \emph{Irrationality and linear independence of Dirichlet beta
values}, GitHub repository, \url{https://github.com/twinheterodox/beta4}.
\end{thebibliography}
\end{document}
